% ============================================================
%  TESIS — UANL / FIME
%  Abraham Castaneda Quintero
%  Muestreo Híbrido mediante Redes de Flujo Generativo
%  y Optimización por Colonia de Hormigas para el Problema
%  de Enrutamiento de Vehículos Estocástico (EHBG-FACS)
%
%  ESQUELETO derivado 100% del diseño de anteproyecto.tex:
%   - mismo preámbulo/portada/colores institucionales;
%   - capítulos mapeados 1:1 a las secciones del anteproyecto
%     (antecedentes → marco teórico; fases 1–5 → caps. 3–6;
%      hipótesis/objetivos/contribuciones → cap. 1 y 7);
%   - el temario de literatura del cap. 2 cubre exactamente los
%     temas y las 14 referencias del anteproyecto.
%  Los huecos por redactar están marcados con \pendiente{...}.
%  Compilar con: pdflatex + biber + pdflatex ×2 (Overleaf: biber).
% ============================================================

\documentclass[12pt]{report}

% ── Paquetes (idénticos al anteproyecto) ────────────────────
\usepackage[spanish,es-noquoting,es-noshorthands]{babel}
\usepackage[utf8]{inputenc}
\usepackage{fontenc}
\usepackage{csquotes}           % requerido por biblatex con babel
\usepackage{graphicx}
\usepackage{geometry}
\usepackage{setspace}
\usepackage{tocloft}
\usepackage[dvipsnames]{xcolor}
\usepackage{tikz}
\usepackage{lmodern}
\usepackage{microtype}
\usepackage{array}
\usepackage{letterspace}
\usepackage{amsmath,amssymb}
\usepackage{enumitem}
\usepackage{booktabs}
\usepackage{multirow}
\usepackage{colortbl}
\usepackage[backend=biber,style=ieee,sorting=none]{biblatex}
\usepackage{hyperref}

% ── Bibliografía (misma base que el anteproyecto) ────────────
\addbibresource{referencias.bib}

% ── Página y tipografía ─────────────────────────────────────
\geometry{letterpaper, margin=1in}
\setlength{\parindent}{1.25cm}
\onehalfspacing
\setlength{\emergencystretch}{3em}

% ── Colores institucionales UANL ────────────────────────────
\definecolor{uanlblue}{RGB}{0,56,101}
\definecolor{uanlgold}{RGB}{200,155,40}
\definecolor{lightbg}{RGB}{245,246,248}
\definecolor{rowgray}{RGB}{235,237,240}

\hypersetup{
  colorlinks = true,
  linkcolor  = uanlblue,
  citecolor  = uanlblue,
  urlcolor   = uanlblue
}

% ── Marcador de secciones por redactar ───────────────────────
\newcommand{\pendiente}[1]{\par\noindent\colorbox{lightbg}{%
  \begin{minipage}{0.94\linewidth}\small
  \textcolor{uanlgold}{\textbf{[POR REDACTAR]}} #1
  \end{minipage}}\par\medskip}

% ════════════════════════════════════════════════════════════
\begin{document}

% ── PORTADA (mismo diseño del anteproyecto) ─────────────────
\begin{titlepage}
\pagecolor{white}

\begin{tikzpicture}[remember picture,overlay]
  \fill[uanlblue] (current page.north west)
      rectangle ++(\paperwidth,-1.2cm);
  \fill[uanlgold] ([yshift=-1.2cm]current page.north west)
      rectangle ++(\paperwidth,-0.22cm);
  \fill[uanlgold] (current page.south west)
      rectangle ++(\paperwidth,0.22cm);
  \fill[uanlblue] ([yshift=0.22cm]current page.south west)
      rectangle ++(\paperwidth,1.2cm);
\end{tikzpicture}

\vspace*{0.5cm}

\begin{center}

{\fontsize{15.5}{19}\selectfont\bfseries\textcolor{uanlblue}{%
  UNIVERSIDAD AUT\'ONOMA DE NUEVO LE\'ON}}\\[0.3cm]
{\fontsize{11.5}{14}\selectfont\textsc{\textcolor{uanlblue}{%
  Facultad de Ingenier\'ia Mec\'anica y El\'ectrica}}}

\vspace{0.6cm}
{\color{uanlgold}\rule{0.55\textwidth}{1.6pt}}
\vspace{0.8cm}

{\fontsize{13.0}{18}\selectfont\bfseries\textcolor{uanlblue}{%
  MUESTREO H\'IBRIDO MEDIANTE REDES DE FLUJO\\[0.2cm]
  GENERATIVO Y OPTIMIZACI\'ON POR COLONIA DE\\[0.2cm]
  HORMIGAS PARA EL PROBLEMA DE ENRUTAMIENTO\\[0.2cm]
  DE VEH\'ICULOS ESTOC\'ASTICO}}

\vspace{0.8cm}
{\color{uanlgold}\rule{0.55\textwidth}{1.6pt}}
\vspace{1.0cm}

{\fontsize{16}{20}\selectfont\bfseries\textcolor{uanlblue}{%
  T\,E\,S\,I\,S}}

\vspace{0.2cm}
{\fontsize{11}{14}\selectfont\textit{Para obtener el t\'itulo de:}}\\[0.45cm]
{\fontsize{13.5}{17}\selectfont\bfseries\textit{%
  Ingeniero en Tecnolog\'ias de Software}}

\vspace{0.5cm}

\colorbox{lightbg}{%
  \begin{minipage}{0.68\textwidth}
    \vspace{0.5cm}
    \centering
    {\small\textit{\textcolor{uanlblue}{Presenta:}}}\\[0.3cm]
    {\fontsize{14}{17}\selectfont\bfseries Abraham Casta\~neda Quintero}\\[0.5cm]
    {\color{uanlgold}\rule{0.7\linewidth}{0.8pt}}\\[0.4cm]
    {\small\textit{\textcolor{uanlblue}{Director de Tesis:}}}\\[0.3cm]
    {\fontsize{12}{15}\selectfont\bfseries Dr. Jos\'e Arturo Berrones Santos}\\[0.5cm]
    \vspace{0.5cm}
  \end{minipage}%
}

\end{center}

\vfill

\begin{center}
  {\small\textcolor{uanlblue}{San Nicol\'as de los Garza, Nuevo Le\'on%
  \enspace\textbullet\enspace 2026}}
\end{center}

\vspace{1.4cm}
\end{titlepage}

% ════════════════════════════════════════════════════════════
% PRELIMINARES (paginación romana)
% ════════════════════════════════════════════════════════════
\pagenumbering{roman}

% ── Hoja de aprobación ───────────────────────────────────────
\chapter*{Hoja de Aprobaci\'on}
\addcontentsline{toc}{chapter}{Hoja de Aprobaci\'on}
\pendiente{Insertar la hoja de aprobaci\'on oficial de la FIME con las firmas
del comit\'e (director, revisores) conforme al formato institucional vigente.}

% ── Dedicatoria y agradecimientos ────────────────────────────
\chapter*{Agradecimientos}
\addcontentsline{toc}{chapter}{Agradecimientos}
\pendiente{Agradecimientos (director de tesis, comit\'e, familia, instituciones).}

% ── Resumen / Abstract ───────────────────────────────────────
\chapter*{Resumen}
\addcontentsline{toc}{chapter}{Resumen}
\pendiente{Resumen en espa\~nol (1 p\'agina): problema (fragilidad de la NCO
determinista ante la estocasticidad urbana), propuesta (EHBG-FACS: GFlowNet de
Balance H\'ibrido + muestreador GFACS + recompensa CVaR, extensi\'on ENN),
m\'etodo (benchmark \texttt{SVRPBench}, protocolo homologado con CRN, validaci\'on
estad\'istica) y hallazgos principales con su significancia.}

\chapter*{Abstract}
\addcontentsline{toc}{chapter}{Abstract}
\pendiente{Traducci\'on fiel del resumen al ingl\'es.}

% ── Índices ──────────────────────────────────────────────────
\tableofcontents
\newpage
\listoffigures
\addcontentsline{toc}{chapter}{\'Indice de figuras}
\newpage
\listoftables
\addcontentsline{toc}{chapter}{\'Indice de cuadros}
\newpage

% ── Lista de acrónimos ───────────────────────────────────────
\chapter*{Lista de Acr\'onimos}
\addcontentsline{toc}{chapter}{Lista de Acr\'onimos}
\begin{tabular}{>{\bfseries\color{uanlblue}}p{3.2cm}p{11.5cm}}
VRP        & \textit{Vehicle Routing Problem} (Problema de Enrutamiento de Veh\'iculos)\\
SVRP       & \textit{Stochastic VRP} (VRP Estoc\'astico)\\
CVRP(TW)   & VRP con capacidad (y ventanas de tiempo)\\
ACO        & \textit{Ant Colony Optimization} (Optimizaci\'on por Colonia de Hormigas)\\
NCO        & Optimizaci\'on Combinatoria Neuronal\\
RL         & \textit{Reinforcement Learning} (Aprendizaje por Refuerzo)\\
AM         & \textit{Attention Model}\\
POMO       & \textit{Policy Optimization with Multiple Optima}\\
GFlowNet   & \textit{Generative Flow Network} (Red de Flujo Generativo)\\
TB / DB    & \textit{Trajectory Balance} / \textit{Detailed Balance}\\
HBG        & \textit{Hybrid-Balance GFlowNet}\\
GFACS      & \textit{Generative Flow Ant Colony Sampler}\\
ENN        & \textit{Epistemic Neural Network} (Red Neuronal Epist\'emica)\\
EHBG-FACS  & \textit{Epistemic} HBG-FACS (marco propuesto)\\
CVaR       & \textit{Conditional Value-at-Risk} (Valor en Riesgo Condicional)\\
CRN        & \textit{Common Random Numbers} (N\'umeros Aleatorios Comunes)\\
B\&C       & \textit{Branch-and-Cut} (Ramificaci\'on y Cortes)\\
MTZ        & Formulaci\'on Miller--Tucker--Zemlin\\
\end{tabular}

% ════════════════════════════════════════════════════════════
% CUERPO (paginación arábiga)
% ════════════════════════════════════════════════════════════
\newpage
\pagenumbering{arabic}

% ────────────────────────────────────────────────────────────
\chapter{Introducci\'on}\label{cap:intro}
% Mapea: "Antecedentes" (¶ inicial), "Definición del Problema",
% "Hipótesis", "Objetivos" y "Contribución" del anteproyecto.
% ────────────────────────────────────────────────────────────

\section{Motivaci\'on y contexto}\label{sec:motivacion}
\pendiente{Redactar a partir del primer bloque de los Antecedentes del
anteproyecto: relevancia industrial del VRP (cadenas de suministro, \'ultima
milla), su NP-dureza~\cite{Lenstra1981Complexity} y la incompatibilidad del
modelo determinista con la log\'istica urbana real~\cite{Heakl2025SVRP}.}

\section{Planteamiento del problema}\label{sec:problema}
\pendiente{Trasladar la ``Definici\'on del Problema de Investigaci\'on'': la
segmentaci\'on tecnol\'ogica entre la NCO determinista (eficiente pero fr\'agil
ante desplazamientos distribucionales) y las GFlowNets (ricas en diversidad y
riesgo, sin validaci\'on en entornos log\'isticos estoc\'asticos de alta
fidelidad).}

\section{Pregunta de investigaci\'on}\label{sec:pregunta}
\begin{center}
  \colorbox{lightbg}{%
    \begin{minipage}{0.85\textwidth}
      \vspace{0.4cm}
      \centering
      \textit{\guillemotleft\textquestiondown Es posible mejorar de forma
      estad\'isticamente significativa la robustez y viabilidad t\'actica de
      soluciones para el problema de ruteo estoc\'astico mediante un marco
      h\'ibrido que combine el muestreo adaptativo de una red neuronal
      Hybrid-Balance GFlowNet con el refinamiento poblacional de colonias de
      hormigas estoc\'asticas?\guillemotright}
      \vspace{0.4cm}
    \end{minipage}%
  }
\end{center}
\pendiente{Definir operacionalmente ``robustez'' (tasa de factibilidad, costo
esperado con recurso $E[c+Q]$ y $\mathrm{CVaR}_\alpha$) y ``viabilidad
t\'actica'' (cumplimiento de restricciones bajo $\xi$ con tiempos de inferencia
competitivos).}

\section{Hip\'otesis}\label{sec:hipotesis}

\subsection*{Hip\'otesis general}
La integraci\'on arquitect\'onica de una Red de Flujo Generativo con Balance
H\'ibrido (HBG) y cuantificaci\'on de incertidumbre epist\'emica (ENN) como
parametrizador a priori dentro de un enjambre estoc\'astico paralelo (ACO),
entrenada bajo la infraestructura desacoplada de \texttt{RL4CO}, producir\'a un
motor de b\'usqueda combinatorio neuronal capaz de superar estad\'isticamente en
robustez (tasa de factibilidad y costo esperado bajo ruido log-normal y
disrupciones de Poisson) a los m\'etodos neuronales deterministas (POMO, AM) y a
las heur\'isticas industriales convencionales en el marco de simulaci\'on
\texttt{SVRPBench}, manteniendo tiempos de inferencia competitivos frente a
dichas heur\'isticas.

\subsection*{Hip\'otesis espec\'ificas}
\begin{enumerate}[label=\textbf{H\arabic*.},leftmargin=1.5cm,itemsep=0.5em]
  \item La incorporaci\'on expl\'icita de la penalizaci\'on de Balance Detallado
  (DB) junto al Balance de Trayectoria (TB) mitigar\'a los picos de p\'erdida
  asint\'oticos en espacios complejos, mejorando la convergencia en instancias
  con ventanas de tiempo r\'igidas.
  \item La retroalimentaci\'on as\'incrona de trayectorias estoc\'asticas desde
  el motor de Colonia de Hormigas hacia el b\'ufer fuera de pol\'itica de la
  GFlowNet incrementar\'a la diversidad del hiperespacio de soluciones generado,
  limitando el colapso de modo.
  \item La cuantificaci\'on activa de la incertidumbre epist\'emica mediante
  Redes Neuronales Epist\'emicas guiar\'a el muestreo hacia regiones poco
  exploradas del espacio de soluciones, incrementando la diversidad de rutas
  candidatas factibles y, en sinergia con el refinamiento poblacional del ACO,
  facilitando la identificaci\'on de rutas de recurso resilientes frente a los
  retrasos log-normales y las disrupciones de Poisson.
\end{enumerate}
\pendiente{A\~nadir la m\'etrica operacional de cada hip\'otesis (H1:
estabilidad/percentiles de la p\'erdida y \'epocas-a-umbral; H2: m\'etrica
expl\'icita de diversidad de las top-$k$ rutas; H3: correlaci\'on entre se\~nal
epist\'emica y cobertura de subgrafos), conforme a la hoja de ruta de
revisi\'on del anteproyecto.}

\section{Objetivos}\label{sec:objetivos}

\subsection*{Objetivo general}
Desarrollar y validar un marco algor\'itmico h\'ibrido que acople una Red de
Flujo Generativo de Balance H\'ibrido (HBG) con un muestreador de Colonia de
Hormigas (\textit{Generative Flow Ant Colony Sampler}, GFACS) ---n\'ucleo
denominado HBG-FACS--- para la optimizaci\'on robusta del Problema de
Enrutamiento de Veh\'iculos Estoc\'astico de escala peque\~na y media, simulado
en alta fidelidad mediante \texttt{SVRPBench}. Como extensi\'on exploratoria se
incorporar\'a un m\'odulo de incertidumbre epist\'emica (ENN), configurando el
marco completo \textit{Epistemic} HBG-FACS (EHBG-FACS), y se evaluar\'a la
escalabilidad a instancias de mayor tama\~no.

\subsection*{Objetivos espec\'ificos}
\begin{enumerate}[label=\textbf{OE\arabic*.},leftmargin=1.8cm,itemsep=0.5em]
  \item Integrar el marco \texttt{SVRPBench} dentro de la librer\'ia modular
  \texttt{RL4CO}, automatizando la generaci\'on de instancias urbanas din\'amicas
  de escala peque\~na y media (50 a 300 nodos).
  \item Dise\~nar e implementar en PyTorch una arquitectura
  codificador-decodificador basada en Transformers, parametrizada como una
  GFlowNet de Balance H\'ibrido (HBG).
  \item Acoplar la pol\'itica estoc\'astica de la GFlowNet con el muestreador
  metaheur\'istico de Colonia de Hormigas (GFACS) y un b\'ufer fuera de
  pol\'itica, conformando el n\'ucleo HBG-FACS.
  \item Cuantificar emp\'iricamente el desempe\~no del modelo en calidad de
  soluci\'on (brecha de optimalidad y costo esperado), tasa de factibilidad y
  tiempos de inferencia, frente a m\'etodos del estado del arte (POMO, AM y
  heur\'isticas cl\'asicas), con validaci\'on estad\'istica.
  \item \textit{(Exploratorio)} Incorporar el m\'odulo epist\'emico (ENN) para
  guiar la exploraci\'on y evaluar la escalabilidad del marco a configuraciones
  multi-dep\'osito y de gran escala ($>500$ nodos), seg\'un la disponibilidad de
  tiempo y c\'omputo.
\end{enumerate}

\section{Contribuciones}\label{sec:contribuciones}
\begin{enumerate}[label=\textbf{C\arabic*.},leftmargin=1.5cm,itemsep=0.5em]
  \item \textbf{Fundamentaci\'on de un Paradigma H\'ibrido Estoc\'astico:}
  formalizaci\'on documentada del acoplamiento HBG--GFACS para el ruteo
  estoc\'astico de alta fidelidad y su unificaci\'on con redes epist\'emicas
  (ENN)~\cite{Zhang2025Hybrid, ENN_GFN_2025}.
  \item \textbf{Integraci\'on As\'incrona de Redes Profundas y
  Metaheur\'isticas:} validaci\'on emp\'irica del uso de distribuciones
  neuronales probabil\'isticas completas para inicializar b\'usquedas
  descentralizadas masivamente paralelas (ACO) en entornos de tr\'afico
  din\'amicos.
  \item \textbf{Liberaci\'on de un Framework Reproducible de Ruteo Tolerante a
  Fallos:} ecosistema \texttt{EHBG-FACS} sobre la estructura desacoplada de
  \texttt{RL4CO}, depositado en repositorio de c\'odigo
  abierto~\cite{Berto2023RL4CO}.
\end{enumerate}

\section{Alcance y limitaciones}\label{sec:alcance}
\pendiente{Declarar: dep\'osito \'unico; escalas comprometidas 50--300 nodos
(multi-dep\'osito y $>500$ nodos solo exploratorio); ventanas de tiempo tratadas
como restricciones \emph{soft} v\'ia recurso de segunda etapa (caveat de escala
de \texttt{SVRPBench}); l\'imites de licencia de Gurobi para los baselines
exactos; presupuesto de c\'omputo (GPU Colab).}

\section{Estructura de la tesis}\label{sec:estructura}
\pendiente{P\'arrafo-gu\'ia: cap.~\ref{cap:marco} marco te\'orico;
cap.~\ref{cap:metodologia} formulaci\'on y metodolog\'ia;
cap.~\ref{cap:experimental} dise\~no experimental; cap.~\ref{cap:resultados}
resultados; cap.~\ref{cap:enn} extensi\'on epist\'emica;
cap.~\ref{cap:conclusiones} conclusiones.}

% ────────────────────────────────────────────────────────────
\chapter{Marco Te\'orico y Estado del Arte}\label{cap:marco}
% TEMARIO DE LITERATURA — cobertura derivada 100% del anteproyecto:
% cada sección corresponde a un bloque de los "Antecedentes", "La
% Barrera de la Estocasticidad" y "El Surgimiento de las GFlowNets".
% ────────────────────────────────────────────────────────────

\section{El Problema de Enrutamiento de Veh\'iculos}\label{sec:vrp}
\pendiente{Cubrir: definici\'on formal del VRP y variantes relevantes (CVRP,
VRPTW); restricciones de capacidad, ventanas de tiempo y jornada; complejidad
computacional y pertenencia a NP-duro~\cite{Lenstra1981Complexity}; panorama de
los tres paradigmas hist\'oricos de soluci\'on (Cuadro~\ref{tab:paradigmas}).}

\begin{table}[ht]
\centering
\caption{Paradigmas de Optimizaci\'on para el Enrutamiento de Veh\'iculos}
\label{tab:paradigmas}
\renewcommand{\arraystretch}{1.4}
\setlength{\tabcolsep}{5pt}
\small
\begin{tabular}{>{\bfseries\color{uanlblue}}p{3.5cm}
                >{\raggedright\arraybackslash}p{4.2cm}
                >{\raggedright\arraybackslash}p{4.2cm}
                >{\centering\arraybackslash}p{3.1cm}}
\toprule
\rowcolor{uanlblue}
{\color{white}\textbf{Paradigma}} &
{\color{white}\textbf{Ventaja Principal}} &
{\color{white}\textbf{Desventaja Cr\'itica}} &
{\color{white}\textbf{Escalabilidad}} \\
\midrule
M\'etodos Exactos (Branch \& Cut) & Garantizan el \'optimo global con certeza matem\'atica. & Intratabilidad temporal por explosi\'on combinatoria. & Muy Baja ($< 100$ nodos) \\
\rowcolor{rowgray}
Metaheur\'isticas Cl\'asicas (ACO, Tabu) & Soluciones de alta calidad sin entrenamiento previo. & Requieren ejecuci\'on desde cero; dise\~no manual. & Moderada ($100$ - $1000$ nodos) \\
NCO supervisada & Inferencia r\'apida tras el entrenamiento. & Dependencia total de datos etiquetados caros. & Baja (Limitada por datos) \\
\rowcolor{rowgray}
NCO por RL determinista (POMO, AM) & Inferencia en milisegundos; entrenamiento sin etiquetas. & Fragilidad extrema ante estocasticidad; colapso de modos. & Alta ($> 1000$ nodos en entornos est\'aticos) \\
\bottomrule
\end{tabular}
\end{table}

\section{M\'etodos exactos}\label{sec:exactos}
\pendiente{Cubrir: ramificaci\'on y acotaci\'on (\textit{branch-and-bound});
planos de corte y \textit{branch-and-cut}; desigualdades de capacidad
redondeada/eliminaci\'on de subtours (RCI/DFJ) como restricciones perezosas;
formulaciones de flujo (no dirigida de dos \'indices) y MTZ con ventanas;
programaci\'on din\'amica; el l\'imite pr\'actico de escala ($<100$ nodos) que
motiva el resto de paradigmas.}

\section{Metaheur\'isticas cl\'asicas}\label{sec:metaheuristicas}
\pendiente{Cubrir: Recocido Simulado, B\'usqueda Tab\'u, Algoritmos Gen\'eticos
(panor\'amica); \textbf{ACO en detalle}~\cite{Dorigo1996AntSystem}: regla de
transici\'on $\tau^{\alpha}\cdot\eta^{\beta}$, evaporaci\'on y dep\'osito de
feromona $Q/L$, Ant System y variantes; limitaciones estructurales: dise\~no
manual de reglas, calibraci\'on artesanal de hiperpar\'ametros y ejecuci\'on
desde cero por instancia (imposibilidad de amortizaci\'on).}

\section{Optimizaci\'on Combinatoria Neuronal}\label{sec:nco}
\pendiente{Cubrir: motivaci\'on de la NCO (reemplazar heur\'isticas manuales por
redes parametrizadas); \textit{Pointer Networks} y aprendizaje supervisado con
su cuello de botella de etiquetas; el giro al RL (REINFORCE, actor-cr\'itico);
el \textit{Attention Model} codificador-decodificador~\cite{Kool2019Attention};
\textbf{POMO}~\cite{Kwon2020POMO}: explotaci\'on de simetr\'ias, m\'ultiples
trayectorias con inicios distintos y l\'inea base compartida de bajo
sesgo/varianza; el ecosistema \texttt{RL4CO} como infraestructura experimental
est\'andar~\cite{Berto2023RL4CO}.}

\section{Ruteo estoc\'astico y medidas de riesgo}\label{sec:svrp}
\pendiente{Cubrir: formulaci\'on del SVRP (velocidades, tiempos de servicio y
demandas como variables aleatorias); \textbf{programaci\'on estoc\'astica en dos
etapas con recurso}~\cite{Gendreau1996Stochastic}: ruta \emph{a priori},
escenario $\xi$, pol\'itica de recurso y costo $c+Q$; medidas de riesgo para
colas pesadas: VaR y \textbf{CVaR}~\cite{Rockafellar2000CVaR}, y por qu\'e el
costo esperado es insuficiente ante disrupciones de cola.}

\section{El marco \texttt{SVRPBench}}\label{sec:svrpbench}
\pendiente{Cubrir~\cite{Heakl2025SVRP}: los cuatro vectores estoc\'asticos
(congesti\'on por mezclas gaussianas dependiente del tiempo; retrasos
log-normales de cola pesada; disrupciones por procesos de Poisson; ventanas
heterog\'eneas residencial/comercial); el hallazgo central: degradaci\'on
$>20\%$ de los modelos NCO deterministas ante desplazamientos distribucionales
(precisar m\'etrica y condici\'on del resultado), y el colapso de modos como
mecanismo; baselines incluidos en el benchmark (OR-Tools, LKH, metaheur\'isticas
oficiales).}

\section{Redes de Flujo Generativo}\label{sec:gfn}
\pendiente{Cubrir: fundamentos de GFlowNets~\cite{Bengio2023Foundations}:
DAG de construcci\'on $\mathcal{G}=(\mathcal{V},\mathcal{E})$, estados
parciales/terminales, conservaci\'on de flujo, pol\'iticas $P_F$/$P_B$ y la
propiedad de muestreo $P_T(x)\propto R(x)$ (en el punto fijo del objetivo);
diferencia conceptual con RL de maximizaci\'on.}

Para entrenar estas pol\'iticas, el objetivo de Balance de Trayectoria
(TB)~\cite{Malkin2022Trajectory} optimiza a lo largo de una trayectoria completa
$\tau = (s_0 \to \dots \to s_n)$:
\begin{equation}
\mathcal{L}_{TB}(\tau; \theta) = \left( \log \frac{Z_\theta \prod_{t=0}^{n-1} P_F(s_{t+1}|s_t; \theta)}{\prod_{t=0}^{n-1} P_B(s_t|s_{t+1}; \theta) R(s_n)} \right)^2
\label{eq:tb}
\end{equation}
mientras que el principio de Balance Detallado (DB) opera a nivel de
transici\'on at\'omica:
\begin{equation}
\mathcal{L}_{DB}(s, s'; \theta) = \left( \log \frac{F_\theta(s) P_F(s'|s; \theta)}{F_\theta(s') P_B(s|s'; \theta)} \right)^2
\label{eq:db}
\end{equation}

\subsection{Balance H\'ibrido (HBG)}\label{sec:hbg}
\pendiente{Cubrir~\cite{Zhang2025Hybrid}: inestabilidades locales del TB
hol\'istico en ruteo; reintroducci\'on del DB y su combinaci\'on adaptativa
(ponderador $\lambda_{DB}$); resultados reportados en VRP.}

\subsection{Muestreo con colonias de hormigas (GFACS)}\label{sec:gfacs}
\pendiente{Cubrir~\cite{Kim2024GFACS}: GFlowNet como generador de la matriz
heur\'istica a priori $\eta$ de un ACO; refinamiento poblacional; posici\'on
frente a DeepACO/heur\'isticas neuronales de aristas.}

\subsection{Incertidumbre epist\'emica (ENN)}\label{sec:enn}
\pendiente{Cubrir~\cite{ENN_GFN_2025}: redes neuronales epist\'emicas y el
esquema \textit{epinet} (cabezas indexadas con cuerpo compartido); su uso para
guiar exploraci\'on en GFlowNets hacia regiones poco visitadas.}

\section{S\'intesis y brecha de investigaci\'on}\label{sec:brecha}
\pendiente{Cerrar el cap\'itulo con la segmentaci\'on tecnol\'ogica del
anteproyecto: NCO determinista = eficiencia sin robustez; GFlowNets = diversidad
y riesgo sin validaci\'on log\'istica de alta fidelidad; el nicho vac\'io que
EHBG-FACS ocupa y c\'omo cada componente de la propuesta responde a una
limitaci\'on documentada en las secciones previas.}

% ────────────────────────────────────────────────────────────
\chapter{Formulaci\'on del Problema y Metodolog\'ia}\label{cap:metodologia}
% Mapea las Fases 1–3 del anteproyecto.
% ────────────────────────────────────────────────────────────

\section{Formulaci\'on: SVRP en dos etapas con recurso}\label{sec:formulacion}
\pendiente{Formalizar (Fase 1): decisi\'on \emph{a priori}; escenario
$\xi$ (congesti\'on, retrasos log-normales, accidentes de Poisson); pol\'itica
de recurso y costo total $c(\text{ruta}) + Q(\text{ruta},\xi)$ estimado por
Monte Carlo~\cite{Gendreau1996Stochastic}; recompensa sensible al riesgo
$R(x)\propto\exp(-\mathrm{CVaR}_\alpha(c+Q)/T)$~\cite{Rockafellar2000CVaR}.}

\section{Ecosistema experimental}\label{sec:ecosistema}
\pendiente{Describir (Fase 1): \texttt{SVRPBench}~\cite{Heakl2025SVRP} sobre
\texttt{RL4CO}~\cite{Berto2023RL4CO}; banco can\'onico de instancias con
semillas fijas (dep\'osito \'unico, escalas 50--300); generaci\'on de los cuatro
vectores estoc\'asticos; n\'umero de realizaciones Monte Carlo por instancia y
uso de N\'umeros Aleatorios Comunes (CRN) para comparabilidad.}

\section{Arquitectura neuronal HBG}\label{sec:arquitectura}
\pendiente{Describir (Fase 2): codificador \textit{Graph Transformer} con
caracter\'isticas est\'aticas y de varianza temporal; decodificador que
parametriza $P_F$, $P_B$, el flujo $F_\theta(s)$ y $\log Z$; p\'erdida h\'ibrida
$\mathcal{L}_{HBG} = (1-\lambda_{DB})\,\mathcal{L}_{TB} +
\lambda_{DB}\,\mathcal{L}_{DB}$ (ecs.~\ref{eq:tb} y~\ref{eq:db})
~\cite{Zhang2025Hybrid}; detalles de entrenamiento (AdamW, precisi\'on mixta).}

\section{Acoplamiento GFACS y b\'ufer fuera de pol\'itica}\label{sec:acoplamiento}
\pendiente{Describir (Fase 3): siembra de la matriz heur\'istica $\eta$ con la
pol\'itica $P_F$~\cite{Kim2024GFACS}; enjambres con iteraciones estoc\'asticas;
actualizaci\'on de feromona por robustez (CVaR); b\'ufer de repetici\'on
off-policy que retroalimenta trayectorias \'elite al entrenamiento HBG;
pseudoc\'odigo del n\'ucleo HBG-FACS.}

\section{Extensi\'on epist\'emica (ENN)}\label{sec:met-enn}
\pendiente{Describir (Fase 5, dise\~no): transformaci\'on de la capa final en
\textit{epinet} con cabezas indexadas~\cite{ENN_GFN_2025} y el mecanismo por el
que la incertidumbre epist\'emica gu\'ia el muestreo. La evaluaci\'on emp\'irica
se reporta en el cap.~\ref{cap:enn}.}

% ────────────────────────────────────────────────────────────
\chapter{Dise\~no Experimental}\label{cap:experimental}
% Mapea la Fase 4 del anteproyecto.
% ────────────────────────────────────────────────────────────

\section{Protocolo homologado}\label{sec:protocolo}
\pendiente{Condiciones id\'enticas para todos los m\'etodos: mismas instancias
(semillas fijas), mismos escenarios $\xi$ (CRN), misma re-puntuaci\'on con el
evaluador estoc\'astico compartido, misma pol\'itica de recurso y mismo costo de
flota; declarar presupuestos de entrenamiento/inferencia por m\'etodo.}

\section{L\'ineas base}\label{sec:baselines}
\pendiente{Describir cada baseline y su rol: exactos Branch \& Cut (cota de
costo y variante con ventanas soft); metaheur\'isticas oficiales ACO y
Tab\'u~\cite{Dorigo1996AntSystem}; NCO supervisada (Attention Model por
imitaci\'on); NCO por RL determinista POMO+AM~\cite{Kwon2020POMO,
Kool2019Attention}; heur\'isticas del benchmark~\cite{Heakl2025SVRP}.}

\section{M\'etricas de evaluaci\'on}\label{sec:metricas}
\pendiente{Definir: brecha de optimalidad y costo determinista; $E[c]$,
$E[c+Q]$, $\mathrm{CVaR}_\alpha$; tasa de factibilidad y tasa de violaci\'on;
robustez (dispersi\'on entre realizaciones); tiempos de inferencia y costo de
entrenamiento amortizado; tama\~no de flota.}

\section{Sintonizaci\'on de hiperpar\'ametros}\label{sec:optuna}
Los hiperpar\'ametros cr\'iticos se sintonizan con Optuna~\cite{Akiba2019Optuna}
sobre el espacio del Cuadro~\ref{tab:espacio}.

\begin{table}[ht]
\centering
\caption{Espacio de b\'usqueda de hiperpar\'ametros del sistema EHBG-FACS}
\label{tab:espacio}
\renewcommand{\arraystretch}{1.4}
\setlength{\tabcolsep}{6pt}
\small
\begin{tabular}{>{\bfseries\color{uanlblue}}p{4.6cm}
                >{\raggedright\arraybackslash}p{4.2cm}
                >{\centering\arraybackslash}p{3.2cm}}
\toprule
\rowcolor{uanlblue}
{\color{white}\textbf{Hiperpar\'ametro Sist\'emico}} &
{\color{white}\textbf{Espacio de B\'usqueda}} &
{\color{white}\textbf{M\'odulo Impactado}} \\
\midrule
Coeficiente HBG ($\lambda_{DB}$) & Uniforme: $[0.1,\,0.9]$ & Ponderaci\'on de p\'erdida TB vs. DB \\
\rowcolor{rowgray}
Tasa de Evaporaci\'on ($\rho$) & Log-Uniforme: $[10^{-3},\,10^{-1}]$ & Din\'amica de feromonas en ACO \\
Temperatura de Flujo ($T$) & Uniforme: $[0.5,\,5.0]$ & Suavizado de la pol\'itica $P_F$ \\
\rowcolor{rowgray}
Tasa de Aprendizaje AdamW & Log-Uniforme: $[10^{-5},\,10^{-3}]$ & Codificador-decodificador Transformer \\
\bottomrule
\end{tabular}
\end{table}

\section{An\'alisis estad\'istico}\label{sec:estadistica}
\pendiente{Plan de la Fase 4: verificaci\'on de supuestos (Shapiro--Wilk,
Levene); ANOVA con $\alpha=0.05$ cuando se cumplen; Friedman + post-hoc de
Wilcoxon pareado con correcci\'on de Holm cuando no (colas pesadas); dise\~no de
bloques por instancia (v\'alido gracias al CRN); n\'umero de instancias por
tama\~no ($\geq 30$) y reporte de tama\~nos del efecto.}

\section{Estudios de ablaci\'on}\label{sec:ablaciones}
\pendiente{Dise\~no de atribuci\'on (hoja de ruta de revisi\'on del
anteproyecto): (a) AM/POMO con recompensa penalizada por ventanas (a\'isla el
efecto del objetivo); (b) GFlowNet con TB puro vs.\ HBG (a\'isla H1); (c)
HBG sin refinamiento GFACS (a\'isla H2); (d) ACO guiado por CVaR sin prior
neuronal (a\'isla el aporte de la red); (e) EHBG-FACS con y sin ENN (a\'isla
H3).}

% ────────────────────────────────────────────────────────────
\chapter{Resultados y Discusi\'on}\label{cap:resultados}
% ────────────────────────────────────────────────────────────

\section{Calidad de soluci\'on y costo esperado}\label{sec:res-costo}
\pendiente{Tablas/figuras por tama\~no: brecha vs.\ exacto, $E[c]$, $E[c+Q]$.}

\section{Robustez ante la estocasticidad}\label{sec:res-robustez}
\pendiente{$\mathrm{CVaR}_\alpha$, tasa de factibilidad, violaciones; el
\emph{tradeoff} costo--factibilidad--flota y la regi\'on ideal.}

\section{Tiempos de c\'omputo}\label{sec:res-tiempos}
\pendiente{Inferencia por m\'etodo y costo de entrenamiento amortizado;
comparativa frente a heur\'isticas (hip\'otesis general, cl\'ausula de
tiempos).}

\section{Contraste de hip\'otesis}\label{sec:res-hipotesis}
\pendiente{Resultados de las pruebas estad\'isticas y de las ablaciones,
organizados por H1, H2, H3 y la hip\'otesis general; efectos y significancia.}

\section{Discusi\'on}\label{sec:discusion}
\pendiente{Lectura frente a la literatura del cap.~\ref{cap:marco}: ¿confirman
los resultados la degradaci\'on de la NCO determinista~\cite{Heakl2025SVRP}?
¿aporta el balance h\'ibrido la estabilidad reportada por~\cite{Zhang2025Hybrid}?
¿reproduce GFACS el beneficio de~\cite{Kim2024GFACS} bajo estocasticidad?
Limitaciones y amenazas a la validez.}

% ────────────────────────────────────────────────────────────
\chapter{Extensi\'on Epist\'emica y Escalabilidad}\label{cap:enn}
% Mapea la Fase 5 (exploratoria) del anteproyecto.
% ────────────────────────────────────────────────────────────

\section{EHBG-FACS con ENN}\label{sec:enn-resultados}
\pendiente{Comparativa HBG-FACS vs.\ EHBG-FACS-ENN: diversidad explorada,
factibilidad y CVaR (H3)~\cite{ENN_GFN_2025}.}

\section{Escalabilidad}\label{sec:escalabilidad}
\pendiente{Comportamiento a $>500$ nodos y configuraciones multi-dep\'osito;
costo computacional; l\'imites observados.}

% ────────────────────────────────────────────────────────────
\chapter{Conclusiones y Trabajo Futuro}\label{cap:conclusiones}
% ────────────────────────────────────────────────────────────

\section{Conclusiones}\label{sec:conclusiones}
\pendiente{Una conclusi\'on por objetivo espec\'ifico (OE1--OE5) y el veredicto
sobre la pregunta de investigaci\'on y la hip\'otesis general.}

\section{Contribuciones alcanzadas}\label{sec:contrib-finales}
\pendiente{Estado final de C1--C3 (formalizaci\'on, validaci\'on emp\'irica,
repositorio liberado).}

\section{Trabajo futuro}\label{sec:futuro}
\pendiente{L\'ineas derivadas: multi-dep\'osito y gran escala como programa
completo; recurso de segunda etapa aprendido; despliegue en datos reales de
tr\'afico.}

% ────────────────────────────────────────────────────────────
\appendix

\chapter{Marco Reproducible}\label{ap:reproducible}
\pendiente{Estructura del repositorio (\texttt{experiments/colab}: notebooks
00--06, paquete \texttt{svrplab}); instrucciones de ejecuci\'on en Colab/GPU;
semillas y protocolo homologado (C3).}

\chapter{Detalles del Banco de Instancias}\label{ap:instancias}
\pendiente{Par\'ametros de generaci\'on (tama\~nos, demandas, capacidad,
ventanas residencial/comercial); semillas por instancia; formato de
persistencia.}

\chapter{Configuraci\'on Experimental Completa}\label{ap:config}
\pendiente{Hiperpar\'ametros finales por m\'etodo; hardware utilizado (GPU/vCPU);
versiones de software; licencias (Gurobi acad\'emica WLS).}

% ────────────────────────────────────────────────────────────
\printbibliography[title={Bibliograf\'ia}]
\addcontentsline{toc}{chapter}{Bibliograf\'ia}

\end{document}
